% ============================================================================
%  6. CONCLUSIONES
%  ----------------------------------------------------------------------------
%  Cierra el informe cubriendo las cuatro actividades del proyecto.
% ============================================================================

\section{Conclusiones}

\subsection{Resultados obtenidos}

El objetivo central del trabajo se cumpli\'o: se implement\'o y valid\'o una
arquitectura de red empresarial funcional y segura sobre GNU/Linux, en la
que los servicios operan de forma integrada y no como piezas aisladas.

En el \'ambito del correo electr\'onico, el servicio qued\'o operativo con
Postfix como agente de transferencia y Dovecot como servidor de acceso a
buzones, con autenticaci\'on \textsc{sasl} y protecci\'on \textsc{starttls}.
El env\'io autenticado, la recepci\'n en formato \texttt{maildir} y el acceso
desde cliente de escritorio se verificaron satisfactory. En el \'ambito de
archivos, los tres protocolos operan sobre un mismo directorio compartido,
de modo que un archivo creado a trav\'es de \textsc{smb}, \textsc{nfs} o
\textsc{http} es inmediatamente accesible desde los otros, lo que evita la
duplicaci\'on de la informaci\'on almacenada.

La asignaci\'on din\'amica de direcciones qued\'o centralizada en el servidor
\textsc{kea}, tras resolverse el conflicto con el servicio \textsc{dhcp} del
router. Los clientes reciben su direcci\'on, la puerta de enlace y el
servidor \textsc{dns} sin intervenci\'on manual, y las direcciones de
infraestructura permanecen estables mediante reservas. El servidor web con
proxy inverso public\'o los servicios internos por nombre, con
redirecci\'on autom\'atica de \textsc{http} a \textsc{https} y una autoridad
de certificaci\'on interna que cubre todos los subdominios.

El despliegue del controlador de dominio Samba 4 fue el hito que permiti\'o
consolidar el resto. Centraliz\'o la identidad, la autenticaci\'on
\textsc{kerberos} y el \textsc{ldap}, y aport\'o el \textsc{dns} autoritativo
del dominio con sus registros de descubrimiento \textsc{srv}. Gracias a
ello, un equipo cliente se une al dominio, resuelve los nombres internos y
accede a los servicios publicados, todo ello a trav\'es de una \'unica
puerta de entrada cifrada, sin configurar direcciones ni puertos en el
cliente.

\subsection{Importancia del proyecto en la formaci\'on profesional}

El proyecto result\'o especialmente formativo porque no se limit\'o a la
configuraci\'on aislada de herramientas, sino que exigi\'o comprender c\'omo
se relacionan los servicios entre s\'i. Un error en la configuraci\'on del
servidor de nombres se manifiesta como un fallo aparentemente ajeno, como
la imposibilidad de autenticar un usuario en el correo; un par\'ametro
mal definido en \textsc{dhcp} impide a los clientes alcanzar al servidor que
precisamente les fue entregado. Comprender esas relaciones es lo que
distingue a un administrador de red de quien simplemente ejecuta
instrucciones.

El entorno de trabajo result\'o especialmente valioso por ser
realista. Las restricciones encontradas --un \textsc{dhcp} activo en el
router que compet\'ia con el del proyecto, la necesidad de sincronizar el
reloj para que \textsc{kerberos} aceptara los tickets, la confianza de la
autoridad de certificaci\'on en cada estaci\'on-- no se deducen de la
teor\'ia: se descubren diagnosticando. Ese constituye el principal
aprendizaje del ejercicio profesional.

\subsection{Aprendizajes adquiridos}

Se adquirieron competencias en cuatro \'ambitos. En el dise\~no de redes,
se aprendi\'o a segmentar una red plana en zonas con pol\'iticas de acceso
distintas, y a justificar cada decisi\'on t\'ecnica con argumentos en lugar
de aplicarla por costumbre.

En la administraci\'on de servicios, se aprendi\'o a configurar correo
electr\'onico, compartici\'on de archivos, asignaci\'on din\'amica de
direcciones, publicaci\'on web y un directorio corporativo,
comprendiendo el rol espec\'ifico que desempeña cada uno dentro del conjunto y
las dependencias que los unen. Se comprendi\'o tambi\'en la diferencia
pr\'actica entre ejecutar un servicio de forma nativa y desplegarlo en un
contenedor, y por qu\'e el controlador de dominio deb\'ia ejecutarse de la
primera forma \cite{sambateam2024}.

En el \'ambito de la seguridad, se aprendi\'o a aplicar la pol\'itica de
rechazo por defecto, a restringir el acceso administrativo a rangos de
direcciones determinados, a aislar los servicios sensibles tras un \'unico
proxy y a confiar el tr\'afico mediante \textsc{tls} con una autoridad
interna \cite{stewart2017}.

En el diagnostico, el principal aprendizaje fue el m\'etodo: leer el
registro del servicio, contrastar la configuraci\'on efectiva con la
esperada y verificar con la herramienta de l\'inea de comandos antes de
modificar nada. Esta rutina result\'o m\'as valiosa que el conocimiento
memorizado de una opci\'n concreta, porque se transfiere a cualquier
servicio que no se haya visto antes \cite{salz2020}.

\subsection{Recomendaciones para futuras implementaciones}

A partir de la experiencia desarrollada se formulan las siguientes
recomendaciones.

\begin{itemize}
  \item \textbf{Automatizar el aprovisionamiento.} Aunque la
        instalaci\'on se realiz\'o de forma manual, el car\'acter
        declarativo de \textsc{samba-tool} y de los archivos de configuraci\'on
        en \textsc{json} de \textsc{kea} permite escribir guiones que
        reconstruyan el dominio completo de forma idempotente, lo que
        resulta indispensable ante la necesidad de recuperar un servidor
        tras un fallo.
  \item \textbf{Automatizar la renovaci\'on de certificados.} El
        certificado \textit{wildcard} generado tiene una vigencia
        acotada. Un proceso programado que lo reemita y lo recargue
        evitar\'ia las interrupciones de servicio por expiraci\'on.
  \item \textbf{Centralizar los registros.} La resoluci\'on de nombres
        externa depende hoy de un \textit{forwarder} hacia el router, un
        punto \'unico de falla. Instalar un \textsc{dns} recursivo local o
        configurar un \textsc{dns} de respaldo eliminar\'ia esa
        dependencia.
  \item \textbf{Implementar el monitoreo de forma efectiva.} El servidor de
        monitoreo est\'a desplegado, pero a\'un no se han configurado
        alertas que notifiquen la ca\'ida de un servicio. Conectar las
        alertas con el servidor de correo cerrar\'ia el ciclo de
        notificaci\'on.
  \item \textbf{Definir y probar la copia de seguridad.} El proceso de
        respaldo debe especificar qu\'e se copia, con qu\'e frecuencia y c\'omo
        se verifica la restauraci\'on. Una copia nunca comprobada no
        constituye un respaldo.
  \item \textbf{Documentar cada cambio.} Mantener un registro de las
        modificaciones aplicadas a la configuraci\'on permiti\'o resolver
        los problemas del cap\'itulo anterior. Convertir esa pr\'actica en
        un procedimiento formalizado es el siguiente paso natural.
\end{itemize}
